Large language model (LLM)-based multi-agent systems (MAS) have emerged as a powerful paradigm for complex problem solving, where prompts play a central role in shaping both individual agent behaviors and inter-agent coordination.
Existing MAS prompt optimization methods primarily optimize prompts before deployment and keep them fixed during execution, making them unable to adapt to errors and state changes revealed by an ongoing trajectory.
We introduce \textbf{MASOpt}, a framework for \textbf{online prompt optimization in multi-agent systems} inspired by classical closed-loop control theory.
MASOpt treats each prompt modification as a control action whose effectiveness must be determined from the subsequent system response.
A lightweight meta-MAS continuously diagnoses the current execution state, proposes local prompt repairs, and verifies their actual effects, allowing ineffective updates to be reverted and unresolved problems to trigger further optimization within the same episode.
To reduce online exploration, we further introduce \textbf{Cross-Actuation Relative Advantage (CARA)}, which compares alternative prompt interventions from matched execution states and learns state-conditioned repair preferences from historical trajectories without updating the underlying LLM parameters.
We formulate an evaluation protocol spanning question answering, code generation, mathematical reasoning, and interactive decision-making benchmarks. The protocol is designed to test whether online closed-loop repair improves over strong offline prompt-optimization baselines.
